%
% 22-polynome.tex -- Polynome und algebraische Funktionen
%
% (c) 2026 Prof Dr Andreas Müller
%
\subsection{Polynome, rationale und algebraische Funktionen}
Die einfachsten komplexen Funktionen sich direkt aus den Grundoperationen
in $\mathbb{C}$.

%
% Potenzen
%
\subsubsection{Potenzen}
\index{Potenz}%
Die ganzzahlige Potenzfunktion
\[
f_n
\colon
\mathbb{C}\to\mathbb{C}
:
z\mapsto z^n
\]
mit $n\in\mathbb{N}$ kann man mithilfe das binomische Satzes
\index{binomischer Lehrsatz}%
als
\begin{align*}
z^n &= u(x,y) + iv(x,y)
\\
&=
\sum_{k=0}^n \binom{n}{k} x^k i^{n-k} y^{n-k}
\end{align*}
schreiben, wobei Real- und Imaginärteil
\begin{align*}
\operatorname{Re}z^n
&=
\sum_{l=0}^{\left\lfloor\frac{n}2\right\rfloor}
(-1)^l
\binom{n}{2l}
x^{2l} y^{n-2l}
\\
\operatorname{Im}z^n
&=
\sum_{l=0}^{\left\lfloor\frac{n-1}2\right\rfloor}
(-1)^l
\binom{n}{2l+1}
x^{2l+1} y^{n-2l-1}
\end{align*}
sind.
Noch einfacher ist jedoch die Beschreibung in Polarform.
\index{Polarform}%
Die $n$-te Potenz ist
\[
z^n
=
\bigl(r(\cos\varphi+i\sin\varphi)\bigr)^n
=
r^n(\cos n\varphi+ i\sin n\varphi),
\]
der Betrag wird mit $n$ potenziert, das Argument wird mit $n$
multipliziert.

Die Polardarstellung suggeriert auch, wie die Potenzfunktion $z^q$ 
\index{Potenzfunktion}%
für beliebige rationale oder reelle Exponenten $q$ definiert werden
soll.
Wir setzen
\begin{equation}
z^q
=
r^q (\cos q\varphi + i \sin q\varphi).
\label{buch:komplex:funktionen:eqn:qpotenz}
\end{equation}

Es ist verführerisch zu denken, dass für $q=\frac12$
die Formel \eqref{buch:komplex:funktionen:eqn:qpotenz}
die Quadratwurzel
oder $q=\frac1n$ die $n$-te Wurzel einer komplexen Zahl ergibt.
Wenn aber eine Zahl $w\in\mathbb{C}$ die Eigenschaft $w^n=z$ hat,
dann hat auch $w\varepsilon_k$ mit
\[
\varepsilon_k
=
\cos\frac{2\pi k}{n}
+
i
\sin\frac{2\pi k}{n}
\]
diese Eigenschaft.
Es ist nämlich
\[
\varepsilon_k^n
=
\cos \biggl(n\cdot \frac{2\pi k}{n}\biggr)
+
i
\sin \biggl(n\cdot \frac{2\pi k}{n}\biggr)
=
\cos 2\pi k
+
i
\sin 2\pi k
=
1,
\]
und daher auch
\[
(w\varepsilon_k)^n
=
w^n\varepsilon_k^n
=
w^n
=
z.
\]
Durch Multiplikation mit den sogenannten $n$-ten \emph{Einheitswurzeln}
\index{Einheitswurzel}%
$\varepsilon_k$ erhält man also $n-1$ weitere Wurzeln.
Für allgemeine rationale Exponenten $q\in\mathbb{Q}$ können auf
analoge Weise mit dem Nenner des Exponenten endlich viele weitere
Lösungen gefunden werden.

Für irrationalen Exponenten suchen wir nach weiteren Lösungen in der
Form
\(
w_1 = w(\cos\varphi + i\sin\varphi)
\).
Die $q$-Potenz ist
\begin{align*}
w_1^q
&=
w^q(\cos q\varphi + i \sin q\varphi)
=
z(\cos q\varphi + i \sin q\varphi).
\end{align*}
Lösungen werden also gefunden für Winkel, die $q\varphi=2\pi k$ mit
$k\in\mathbb{Z}$ erfüllen, oder
\[
\varphi_k = \frac{2\pi k}{q},\qquad k\in\mathbb{Z}.
\]
Falls $q$ irrational ist, sind alle Winkel $\varphi_k$ verschieden,
es gibt in diesem Falle also unendlich viele verschiedene Lösungen.

%
% Polynome
%
\subsubsection{Polynome}
Polynome sind Linearkombinationen von Potenzfunktionen mit ausschliesslich
natürlichen Exponenten.
Ein Polynom
\[
f(z)
=
f_nz^n + f_{n-1}z^{n-1} + \dots + f_1z + f_0
\in
K[z]
\]
definiert also eine in ganz $\mathbb{C}$ definierte komplexe
Funktion.
\index{Polynom}%

%
% Rationale Funktionen
%
\subsubsection{Rationale Funktionen}
Der Quotient zweier Polynome $p(z),q(z)\in K[z]$  definiert eine Komplexe
Funktion
\[
f(z) = \frac{p(z)}{q(z)}.
\]
Sie ist definiert für alle komplexen Zahlen $z$, für die $q(z)\ne 0$ ist.
Der Definitionsbereich ist also
\[
U = \mathbb{C} \setminus \{ z\in\mathbb{C} \mid q(z)=0\}.
\]

\begin{definition}[rationale Funktion]
Sei $K$ ein Körper, dann heisst der Quotient
\[
f(z) = \frac{p(z)}{q(z)}
\]
zweier Polynome $p(z),q(z)\in K[z]$ \emph{rationale Funktion}.
\index{rationale Funktion}%
\index{Funktion!rational}%
Die Menge der rationalen Funktionen wird mit
\[
K(z)
=
\biggl\{
\frac{p(z)}{q(z)}
\;\bigg|\;
p(z),q(z)\in K[z]
\biggr\}
\]
bezeichnet.
\end{definition}

Durch Ausführen der Polynomdivision mit Rest und durch Kürzen kann 
\index{Polynomdivision}%
eine rationale Funktion $f(z)$ immer eindeutig in die Form
\[
f(z)
=
p(z)
+
\frac{r(z)}{q(z)}
\]
gebracht werden, wobei $\deg r(z)<\deg q(z)$ ist und $q$ ein normiertes
Polynom ist.

Rationale Funktionen sind nur Brüche von Polynomen.
Als solche können Sie addiert, subtrahiert, multipliziert und
auch dividiert werden.
Die Menge $K(z)$ ist also sogar ein Körper.
Die Notation ist gerechtfertigt dadurch, dass man rationale
Funktionen auch als algebraische Ausdrücke betrachten kann,
die dadurch entstehen, dass man den Körper $K$ um die Variable
$z$ erweitert, für die keine weiteren algebraischen Identitäten
gelten.

%
% Algebraische Funktionen
%
\subsubsection{Algebraische Funktionen}
Die Formel~\eqref{buch:komplex:funktionen:eqn:qpotenz}
suggeriert, dass es möglich ist, beliebige Wurzelfunktionen zu definieren.
Die Wurzel $w=f(z)=z^{1/n}$ erfüllt die Polynomgleichung
\[
f(z)^n - z = 0.
\]
Die Definition macht auch klar, dass es jeweils mehrere Wurzeln
gibt.
Immerhin können wir eine Formel dafür angeben, die Lösung der
Polynomgleichung
\[
f(p,q)^2 + pf(p,q) + q = 0
\quad\Rightarrow\quad
f_\pm(p,q)
=
-\frac{p}{2}
\pm
\!\sqrt{\rlap{$\phantom{\bigg|}$}\smash{\biggl(\frac{p}2\biggr)^2 - q}}
\]
schreiben.
Die Lösung ist eine komplexe Funktion der möglicherweise
komplexen Paramter $p$ und $q$.
Falls $p$ und $q$ komplexe Funktionen einer einzelnen
komplexen Variable $z$ sind, dann ist
$f(z) = f_\pm(p(z),q(z))$ eine komplexe Funktion, die die Gleichung
\[
f(z)^2 + p(z) f(z) + q(z) = 0
\]
erfüllt.
Die Lösungsformel verwendet nur die arithmetischen Grundoperationen
und Wurzeln.
So eine Lösung wird eine Lösung durch Radikale genannt.
Eine ähnliche Formel, eine Lösung durch Radikale, ist auch für
die kubische Gleichung und für die Gleichung vierten Grades bekannt.
Niels Henrik Abel hat aber 1824 bewiesen, dass es im Allgemeinen
\index{Abel, Niels Henrik}%
nicht möglich ist, eine Lösung einer Polynomgleichung durch
Radikale anzugeben, wenn der Grad grösser als $4$ ist.

Eine Polynomgleichung höheren Grades der Form
\[
a_n(z) f(z)^n + a_{n-1}(z) f(z)^{n-1} + \dots + a_1(z) f(z) + a_0(z)
=
0,
\]
deren Koeffizienten $a_k(z)$ komplexe Funktionen sind,
definiert eine Lösungsfunktion $f(z)$, die jedoch nicht unbedingt
mit einer einfachen Formel ausgedrückt werden kann.

\begin{definition}[algebraische Funktion]
Eine Funktion $f(z)$ heisst \emph{algebraisch} über einem Körper
\index{algebraische Funktion}%
\index{Funktion!algebraisch}%
$F$ von komplexen Funktionen, wenn es ein Polynom $p(X)\in F[X]$
der Form
\[
p(X)
=
X^n + p_{n-1}(z) X^{n-1} + \dots + p_1(z) X + p_0(z)
\]
gibt, so dass 
\[
p(f(z)) = 0
\]
identisch erfüllt ist.
\end{definition}

Da die $n$-te Wurzel $f(z)$ die Polynomgleichung
\[
p(X)
=
X^n - z
\quad\Rightarrow\quad
p_n(z)=1, p_0(z)=z
\quad\Rightarrow\quad
p(f(z))
=
f(z)^n - z
=
0.
\]
Insbesondere sind die Wurzelfunktionen also algebraisch über dem 
Körper $K(z)$.
